Deep learning has revolutionized photoacoustic image reconstruction from raw data obtained from sensors, where one axis represents time and the other the scanning direction. Approaches have evolved from using conventional methods complemented with deep learning to improve image quality, to developing neural networks that learn to reconstruct images directly from raw data. In this context, the U-Net architecture has become a de facto standard.

However, the main challenge when using neural networks for direct reconstruction lies in the underlying physics of the problem, which makes it mathematically ill-conditioned. Although methods such as Physics Informed Neural Networks (PINN) exist to incorporate physical constraints, these are usually computationally expensive due to the numerous numerical operations required. In this work, we propose a new framework that allows the neural network to not only reconstruct the image but also maintain the physical consistency of the problem.

Our proposal is based on an approach inspired by data guided learning, consisting of two main phases. In the first, we train two U-Net networks: one learns to reconstruct images from raw data, while the other learns the inverse problem, i.e., the generation of photoacoustic signals from images, thus simulating the physical process of acoustic wave propagation. In the second phase, we perform fine-tuning of the reconstruction network, incorporating physical consistency as an additional constraint in the learning process. This approach acts as a supervisor that guides the network to not only learn to reconstruct high-quality images, but also respect the fundamental physical principles of the photoacoustic process.

\textbf{Keywords:} Deep learning, photoacoustic image reconstruction, U-Net, Physics Informed Neural Networks, data guided learning, physical consistency.